% Part: many-valued-logic
% Chapter: syntax-and-semantics
% Section: formulas

\documentclass[../../../include/open-logic-section]{subfiles}

\begin{document}

\olfileid{mvl}{syn}{fml}

\olsection{\usetoken{P}{formula}}

\begin{defn}[সূত্র]
\ollabel{defn:formulas}
কোনো বচনমূলক ভাষা~$\Lang L$-এর \emph{!!{formula}s}-এর সেট~$\Frm[L]$
নিচের মতো আরোহীভাবে সংজ্ঞায়িত:
\begin{enumerate}
\item প্রতিটি !!{propositional variable}~$\Obj p_i$ একটি পরমাণু
  !!{formula}।
\item $\Lang L$-এর প্রতিটি $0$-স্থানীয় সংযোজক (বচনধ্রুবক) একটি পরমাণু
  !!{formula}।
\item $\star$ যদি $\Lang L$-এর একটি $n$-স্থানীয় সংযোজক হয় এবং $!A_1$,
\dots, $!A_n$ যদি !!{formula}s হয়, তবে $\star(!A_1, \dots, !A_n)$
  !!a{formula}।
\tagitem{limitClause}{অন্য কিছুই !!a{formula} নয়।}{}
\end{enumerate}
$\star$ যদি $1$-স্থানীয় হয়, তবে $\star(!A_1)$-কে প্রায়ই সংক্ষেপে
$\star !A_1$ লেখা হবে। $\star$ যদি $2$-স্থানীয় হয়, তবে
$\star(!A_1,!A_2)$-কে প্রায়ই $(!A_1 \star !A_2)$ লেখা হবে।
\end{defn}

প্রচলিত রীতি অনুযায়ী আমরা প্রায়ই সবচেয়ে বাইরের বন্ধনী নিঃশব্দে বাদ দেব।

\begin{ex}
  মানক ভাষা~$\Lang{L_0}$-তে $\Obj p_1 \lif (\Obj p_1
  \land \lnot \Obj p_2)$ একটি সূত্র। গুণন যুক্তিবিদ্যার ভাষায় সেটি
  বরং $\Obj p_1 \lif (\Obj p_1 \odot
  \lnot \Obj p_2)$ লেখা হবে। ভাষায় $1$-স্থানীয়~$\triangle$ যোগ করলে
  $\triangle (\Obj p_1
  \land \Obj p_2) \lif (\triangle \Obj p_1 \land \triangle \Obj p_2)$-এর
  মতো সূত্রও পাওয়া যাবে।
\end{ex}

\end{document}
